% ID: SEL-D0025-OB-E02
% ESTADO: validada. S026, cotejada con examen y criterios.
\textbf{a)} El dominio es $\mathbb R$. La expresión racional
sólo se usa en $x\le0$, de modo que su denominador nunca se
anula en ese tramo. En $0$ el límite izquierdo y el valor son
$-1/4$, mientras que el límite derecho es $3$. Hay una
\textbf{discontinuidad de salto finito}. En $2$, el límite
izquierdo es $2+3=5$ y el derecho y el valor son $2^2+1=5$.
Por tanto es continua en $2$ y en todo $\mathbb R\setminus\{0\}$.

\textbf{b)} En $0$ la discontinuidad impide que exista derivada.
En $2$ las pendientes laterales son $1$ y $4$, por lo que
tampoco es derivable. En el resto del dominio se puede derivar:
$-1/(x-4)^2$ para $x<0$, $1$ para $0<x<2$ y $2x$ para $x>2$.
